% !TeX root = ../../Book2.tex
% 纲要标题：### 16.3 Morita 等价定理
% 纲要位置：计划纲要.md，第 1576 行。

\section{Morita 等价定理}
\label{b2:sec:1603}

第五章把 Morita 等价定义为模范畴的等价，却没有说明怎样从环本身判断它。本节证明：选择一个有限生成投射生成元，再取其自同态环的反环，恰好得到全部这样的等价。先从一个已给的范畴等价恢复这个模，再证明投射生成元确实使上一节的单位、余单位在每个对象上成为同构。

% 纲要要点（供目录核对）：
%
% * 从范畴等价恢复投射生成元；由双积恢复加法，先保持紧性和投射性，再恢复有限生成性
% * 由投射生成元构造等价；将单位与余单位从生成对象经任意直和、展示及余核推广到全部模
% * 互逆双模及右模版本
% * 满幂等元与角环的 Morita 等价；正式纳入非零对象检测判据和求值逆，并证明双边理想及商环对应
%

\subsection{从范畴等价恢复投射生成元}
\label{b2:subsec:1603:01}

\begin{lemma}\label{b2:lem:module-equivalence-structure}
设 \(R,S\) 为含幺环，\(H:R\text{-}\mathbf{Mod}\rightarrow S\text{-}\mathbf{Mod}\) 为幺左模范畴之间的等价，则 \(H\) 为加性函子，保持并反映核、余核、正合列、投射性与内射性，也保持任意直和，并将生成元送到生成元。
\end{lemma}
\begin{proof}
核、余核及直和由泛性质刻画，投射性、内射性与生成元也可用同态的提升、延拓和商对象刻画，因此这些性质都能随等价转移。需要先说明的是，同态的加法本身也能由范畴结构恢复。由\cref{b2:prop:equivalence-preserves-universal-objects}，等价保持零对象及任意积、余积，特别保持任意直和。零同态是经过零对象的复合，因此也被保持。对同态 \(f,g:M\rightarrow N\)，加法可写成
\[
f+g=\nabla_N(f\oplus g)\Delta_M,
\]
其中 \(\Delta_M:M\rightarrow M\oplus M\) 的两个投影均为恒等，\(\nabla_N:N\oplus N\rightarrow N\) 的两个限制均为恒等。直和同时是积与余积，故等价在相应的自然直和识别下保持这两个映射及此式，所以保持加法。无需在模范畴的等价之外另加一个加性要求。

等价保持并反映单态射和满态射的消去性质；模范畴中的单态射就是单射，因为非零核元素可以通过从 \(R\) 出发的同态检测。满态射就是满射，因为若像不为全体，非零商映射与零映射在其像上相同，会破坏右消去性。因此单射与满射也被保持并反映。

以核为例，设 \(j:K\rightarrow M\) 是 \(f:M\rightarrow N\) 的核，等价记为 \(H\)。若 \(a:Y\rightarrow H(M)\) 满足 \(H(f)a=0\)，取同构 \(u:H(X)\rightarrow Y\)，再由全忠实性把 \(au\) 唯一提升为 \(b:X\rightarrow M\)。忠实性给出 \(fb=0\)，故 \(b=jc\) 且 \(c\) 唯一。于是 \(a=H(j)H(c)u^{-1}\)，而任何别的分解提升后都等于 \(c\)。所以 \(H(j)\) 是核。

对余核 \(q:N\rightarrow C\)，给定 \(a:H(N)\rightarrow Y\) 且 \(aH(f)=0\)，选 \(u:H(X)\rightarrow Y\)，把 \(u^{-1}a\) 提升为 \(b:N\rightarrow X\)，则 \(bf=0\)，由余核性质唯一写成 \(b=cq\)，所以 \(a=uH(c)H(q)\) 且分解唯一。故余核也保持。像可表为余核的核，因此也保持。对拟逆使用同样论证，得到反映性，正合列因核与像的相等而双向保持。

投射性由沿满射提升的性质刻画。将一个提升问题经拟逆送回，满射仍满，再作提升并经等价送出，就保持投射性；反向相同。对偶地，内射性由沿单射延拓的性质刻画；等价及拟逆都保持单射，故也保持并反映内射性。若 \(P\) 是生成元，每个源范畴对象都为 \(P\) 某个直和的商。等价保持直和与满射，且每个目标对象都同构于某个 \(H(X)\)，故每个目标对象都为 \(H(P)\) 某个直和的商。因此 \(H(P)\) 也是生成元。
\end{proof}

目标范畴的正则模 \({}_SS\) 是投射生成元，而且 \(\operatorname{Hom}_S(S,Y)\) 在 \(1\) 处求值就恢复 \(Y\)。因此把它经拟逆送回得到 \(P=\Psi(S)\)，应当能用 \(\operatorname{Hom}_R(P,-)\) 恢复原等价。投射性与生成元性质已有范畴刻画；有限生成性则不能直接用元素和生成元的个数转移，需要先把正则模的紧性转移给 \(P\)，再利用 \(P\) 的投射性。

\begin{theorem}[由范畴等价恢复投射生成元]\label{b2:thm:morita-converse}
设 \(R,S\) 为含幺环，\(\Phi:R\text{-}\mathbf{Mod}\rightarrow S\text{-}\mathbf{Mod}\) 为幺左模范畴的等价，取其拟逆 \(\Psi\)，令 \(P=\Psi({}_SS)\)，则幺左 \(R\) 模 \(P\) 为有限生成投射生成元，并有环同构
\[
S\cong\operatorname{End}_R(P)^{\mathrm{op}}.
\]
通过这个同构赋予 Hom 左 \(S\) 作用后，\(\Phi\) 自然同构于 \(\operatorname{Hom}_R(P,-)\)。
\end{theorem}
\begin{proof}
正则模 \(S\) 投射且为生成元，\cref{b2:lem:module-equivalence-structure}使 \(P=\Psi(S)\) 也投射且为生成元。固定拟逆提供的自然同构 \(\nu:\operatorname{Id}_{S\text{-}\mathbf{Mod}}\rightarrow\Phi\Psi\)。为证明有限生成，先证明 \(P\) 紧。对任意幺左 \(R\) 模族 \((M_i)\)，有自然同构
\[
\begin{aligned}
\operatorname{Hom}_R\left(P,\bigoplus_iM_i\right)
&\cong\operatorname{Hom}_S\left(S,\Phi\left(\bigoplus_iM_i\right)\right)\\
&\cong\operatorname{Hom}_S\left(S,\bigoplus_i\Phi(M_i)\right)\\
&\cong\bigoplus_i\operatorname{Hom}_S\bigl(S,\Phi(M_i)\bigr)\\
&\cong\bigoplus_i\operatorname{Hom}_R(P,M_i).
\end{aligned}
\]
第一项同构来自 \(\Phi\) 的全忠实性及 \(\nu_S:S\cong\Phi(P)\)，第二项来自等价保持直和。第三项使用正则模 \(S\) 由 \(1\) 生成：从 \(S\) 到直和的同态由 \(1\) 的像决定，而这个像只有有限支集。最后一项再用全忠实性及 \(\nu_S\)。这些同构与各直和嵌入相容，所以其复合正是紧性中的自然映射之逆。\cref{b2:prop:compact-projective-finite}于是说明投射模 \(P\) 有限生成。

等价加性且全忠实，故 \(\Psi\) 在自同态环上给出环同构
\[
\operatorname{End}_S({}_SS)\cong\operatorname{End}_R(P).
\]
由\cref{b2:prop:regular-hom-endomorphism}，正则左模的自同态是右乘映射 \(r_s:x\mapsto xs\)，并有 \(r_s\circ r_t=r_{ts}\)。因此 \(s\mapsto\Psi(r_s)\) 给出 \(S\cong\operatorname{End}_R(P)^{\mathrm{op}}\)，相应的右作用为 \(ps=\Psi(r_s)(p)\)。

对每个 \(M\)，再用正则模的求值对应
\[
\operatorname{Hom}_R(P,M)
\cong\operatorname{Hom}_S(S,\Phi(M))
\cong\Phi(M).
\]
具体地，\(h:P\rightarrow M\) 先对应 \(a=\Phi(h)\nu_S:S\rightarrow\Phi(M)\)，再对应 \(a(1)\)。后复合 \(M\rightarrow M'\) 的同态与这两步相容，故所得同构对 \(M\) 自然。它也保持左 \(S\) 作用：Hom 中的 \(s\cdot h\) 为 \(h\circ\Psi(r_s)\)，而 \(\nu\) 的自然性给出
\[
\Phi\bigl(h\circ\Psi(r_s)\bigr)\nu_S
=\Phi(h)\nu_S r_s=a r_s.
\]
在 \(1\) 处求值，右边为 \(a(s)=s\cdot a(1)\)，所以该自然同构为左 \(S\) 模同构。
\end{proof}

\subsection{投射生成元与等价构造}
\label{b2:subsec:1603:02}

\begin{theorem}[Morita 等价]\label{b2:thm:morita-equivalence}
设 \(R\) 为含幺环，\(P\) 为幺左 \(R\) 模范畴的有限生成投射生成元，令 \(S=\operatorname{End}_R(P)^{\mathrm{op}}\)，并使 \(s\in S\) 在 \(P\) 上的右作用为 \(ps=\alpha_s(p)\)，其中 \(\alpha_s\) 是 \(s\) 对应的左 \(R\) 线性自同态，则
\[
F=\operatorname{Hom}_R(P,-),\qquad G=P\otimes_S-
\]
分别为从幺左 \(R\) 模到幺左 \(S\) 模、从幺左 \(S\) 模到幺左 \(R\) 模的互逆等价，而且式\eqref{b2:eq:morita-unit}与式\eqref{b2:eq:morita-counit}的单位、余单位在每个模上都是同构。结合\cref{b2:thm:morita-converse}，任意两个含幺环 \(R,S\) 在 Morita 意义下等价，当且仅当存在幺左 \(R\) 模的有限生成投射生成元 \(P\) 使 \(S\cong\operatorname{End}_R(P)^{\mathrm{op}}\)。
\end{theorem}
\begin{proof}
上一节已知 \(\varepsilon_P\) 与 \(\eta_S\) 为同构。要把它们扩展到所有对象，可以先扩展到 \(P\) 或 \(S\) 的任意直和，再把任意模写成这些直和之间某个映射的余核。这里有限生成性与\cref{b2:prop:compact-projective-finite}保证 \(F\) 保持任意直和；投射性与\cref{b2:cor:projective-hom-exact}保证 \(F\) 正合。\cref{b2:prop:tensor-direct-sums}与\cref{b2:thm:tensor-right-exact}说明 \(G\) 保持直和且右正合。因此 \(FG,GF\) 都保持直和与正合尾段。

由\cref{b2:prop:morita-adjunction-maps}，\(\varepsilon_P\) 为同构；因两边函子都保持直和，且求值公式与各嵌入相容，每个 \(\varepsilon_{P^{(I)}}\) 也是同构。对任意幺左 \(R\) 模 \(M\)，生成元性质给出满射 \(q:P^{(J)}\rightarrow M\)。记其核为 \(K\)，再对 \(K\) 使用生成元性质，得到满射 \(P^{(I)}\rightarrow K\)。将它与核的嵌入复合，得到 \(d:P^{(I)}\rightarrow P^{(J)}\) 及正合列
\[
P^{(I)}\longrightarrow P^{(J)}\longrightarrow M\longrightarrow0.
\]
生成元性质的这两次使用，分别保证 \(M\) 有由 \(P\) 给出的生成元和关系。应用右正合函子 \(GF\)，再用余单位的自然性，得到交换的正合尾段
\[
\begin{tikzcd}[column sep=2.6em,row sep=2.5em]
GF(P^{(I)}) \arrow[r,"GF(d)"] \arrow[d,"\varepsilon_{P^{(I)}}"'] &
GF(P^{(J)}) \arrow[r,"GF(q)"] \arrow[d,"\varepsilon_{P^{(J)}}"'] &
GF(M) \arrow[r] \arrow[d,"\varepsilon_M"'] & 0\\
P^{(I)} \arrow[r,"d"'] &
P^{(J)} \arrow[r,"q"'] &
M \arrow[r] & 0
\end{tikzcd}
\]
前两列的竖直映射为同构，并把 \(\operatorname{im}GF(d)\) 对应到 \(\operatorname{im}d\)，因而在对这两个像取商后诱导同构。两行的最后一个对象都是相应第一映射的余核，故由余核泛性质，这个同构正是 \(\varepsilon_M\)。

已知 \(\eta_S\) 为同构，两边保持直和，使每个 \(\eta_{S^{(I)}}\) 也为同构。任意幺左 \(S\) 模 \(N\) 本来就有自由展示
\[
S^{(I)}\longrightarrow S^{(J)}\longrightarrow N\longrightarrow0,
\]
这里指标集允许无限；这一侧用的是正则模 \(S\) 的生成元性质，不再需要 \(P\) 生成 \(R\) 模。把前两个映射记为 \(d':S^{(I)}\to S^{(J)}\)、\(q':S^{(J)}\to N\)。应用右正合函子 \(FG\)，由单位的自然性得到
\[
\begin{tikzcd}[column sep=2.6em,row sep=2.5em]
S^{(I)} \arrow[r,"d'"] \arrow[d,"\eta_{S^{(I)}}"'] &
S^{(J)} \arrow[r,"q'"] \arrow[d,"\eta_{S^{(J)}}"'] &
N \arrow[r] \arrow[d,"\eta_N"'] & 0\\
FG(S^{(I)}) \arrow[r,"FG(d')"'] &
FG(S^{(J)}) \arrow[r,"FG(q')"'] &
FG(N) \arrow[r] & 0
\end{tikzcd}
\]
前两列的竖直映射为同构，第三列是它们在余核上诱导的映射，故 \(\eta_N\) 也为同构。

所以 \(\eta:\operatorname{Id}\rightarrow FG\)、\(\varepsilon:GF\rightarrow\operatorname{Id}\) 都是自然同构，正好给出两函子互为拟逆。任意等价的反方向已经由\cref{b2:thm:morita-converse}证明，得到最后的充要条件。
\end{proof}

这就是\term[Morita-dengjia-dingli]{Morita 等价定理}[Morita equivalence theorem]。仅在正则模或有限自由模上检验单位、余单位，不足以替代证明中通过任意直和和余核所作的扩展。

\ProseParagraphBreak
同态也可显式恢复。若 \(a:F(M)\rightarrow F(M')\) 左 \(S\) 线性，则它唯一来自
\[
\varepsilon_{M'}\,G(a)\,\varepsilon_M^{-1}:M\longrightarrow M'.
\]
确实，式\eqref{b2:eq:morita-triangles}给出 \(F(\varepsilon_M)=\eta_{F(M)}^{-1}\)，单位的自然性便说明将上式施加 \(F\) 后得到 \(a\)。唯一性来自等价的忠实性。这使对象对应与同态对应同时可计算。

\begin{corollary}[互逆双模]\label{b2:cor:morita-inverse-bimodules}
设 \(R\) 为含幺环，\({}_RP\) 为幺左模范畴的有限生成投射生成元，\(S=\operatorname{End}_R(P)^{\mathrm{op}}\)，并令 \(ps=\alpha_s(p)\)，其中 \(\alpha_s\in\operatorname{End}_R(P)\) 对应 \(s\in S\)。令 \(Q=P^*=\operatorname{Hom}_R(P,R)\)，其幺 \((S,R)\) 双模作用为 \((s\lambda)(p)=\lambda(ps)\)、\((\lambda r)(p)=\lambda(p)r\)。则有双模同构
\[
{}_RP\otimes_S Q_R\cong {}_RR_R,\qquad
{}_SQ\otimes_RP_S\cong {}_SS_S.
\]
第一式把 \(p\otimes\lambda\) 送到 \(\lambda(p)\)；第二式把 \(\lambda\otimes p\) 送到 \(S\) 中对应于 \(x\mapsto\lambda(x)p\) 的元素。
\end{corollary}
\begin{proof}
第一式是\eqref{b2:eq:morita-counit}在 \(M=R\) 时的余单位，由\cref{b2:thm:morita-equivalence}为左 \(R\) 模同构。它也保持右 \(R\) 作用，因为 \((\lambda r)(p)=\lambda(p)r\)。

由\cref{b2:prop:finite-projective-hom-tensor}取 \(M=P\)，映射
\[
Q\otimes_RP\longrightarrow\operatorname{End}_R(P),\qquad
\lambda\otimes p\longmapsto T_{\lambda,p},\quad
T_{\lambda,p}(x)=\lambda(x)p
\]
为同构。有限对偶基解释了为什么这些自同态足以构成所有自同态。由\cref{b2:prop:finite-dual-basis}取 \((p_i,\lambda_i)\)，则 \(x=\sum_i\lambda_i(x)p_i\)，因而每个 \(R\) 线性自同态 \(T\) 都满足
\[
T(x)=\sum_i\lambda_i(x)T(p_i).
\]
所以任意自同态都是上述形式的有限和，且已证同构的逆将 \(T\) 送到 \(\sum_i\lambda_i\otimes T(p_i)\)。

将自同态 \(T\) 看作反环 \(S\) 的元素时，左、右 \(S\) 作用分别对应于两种复合次序：
\[
T_{s\lambda,p}=T_{\lambda,p}\circ\alpha_s,\qquad
T_{\lambda,ps}=\alpha_s\circ T_{\lambda,p}.
\]
由于 \(S\) 的乘法与自同态复合反向，这两式分别是 \(S\) 中的左乘 \(s\) 与右乘 \(s\)。故第二式是 \((S,S)\) 双模同构；这里 \(\operatorname{End}_R(P)\) 与 \(S\) 的识别用的是其加性群及这两个作用，而不是把两者的环乘法认作相同。
\end{proof}

这两个同构说明 \({}_RP_S\) 与 \({}_SQ_R\) 在双模张量意义下互为逆：它们的两种复合分别恢复正则双模 \({}_RR_R\) 与 \({}_SS_S\)。同一对双模因此也能从右模一侧构造等价。

\begin{corollary}\label{b2:cor:morita-left-right}
对任意含幺环 \(R,S\)，全部幺左模范畴等价当且仅当全部幺右模范畴等价：
\[
R\text{-}\mathbf{Mod}\simeq S\text{-}\mathbf{Mod}
\quad\Longleftrightarrow\quad
\mathbf{Mod}\text{-}R\simeq\mathbf{Mod}\text{-}S.
\]
\end{corollary}
\begin{proof}
若左模范畴等价，由\cref{b2:thm:morita-converse,b2:cor:morita-inverse-bimodules}取互逆双模 \({}_RP_S\) 与 \({}_SQ_R\)。函子 \(-\otimes_RP\) 把幺右 \(R\) 模送到幺右 \(S\) 模，\(-\otimes_SQ\) 把幺右 \(S\) 模送到幺右 \(R\) 模。两次复合经张量结合律分别同构于 \(-\otimes_R(P\otimes_SQ)\cong-\otimes_RR\) 和 \(-\otimes_S(Q\otimes_RP)\cong-\otimes_SS\)，故互为拟逆。反过来，把右 \(R,S\) 模分别视为左 \(R^{\mathrm{op}},S^{\mathrm{op}}\) 模，对刚证的方向再用一次，便得到左模范畴的等价。
\end{proof}

\subsection{满幂等元与角环的 Morita 等价}
\label{b2:subsec:1603:03}

\begin{proposition}[满幂等角环]\label{b2:prop:full-idempotent-morita}
设 \(R\) 为含幺环，\(e\in R\) 满足 \(e^2=e\)，令 \(S=eRe\)，其单位元为 \(e\)。条件 \(ReR=R\)、幺左 \(R\) 模 \(Re\) 为生成元，以及对每个幺左 \(R\) 模 \(M\) 都有 \(eM=0\Longrightarrow M=0\)，三者等价。在这些条件下，\(Re\) 为有限生成投射生成元，且
\[
R\text{-}\mathbf{Mod}\simeq S\text{-}\mathbf{Mod}
\]
由 \(M\mapsto eM\) 与 \(N\mapsto Re\otimes_SN\) 实现，其中 \(eM\) 取自然的幺左 \(S\) 模结构。对每个幺左 \(R\) 模 \(M\)，求值映射
\[
\varepsilon_M:Re\otimes_SeM\longrightarrow M,\qquad
re\otimes m\longmapsto rm
\]
为自然的左 \(R\) 模同构。
\end{proposition}
\begin{proof}
由 \(R=Re\oplus R(1-e)\)，\(Re\) 是正则左模的直和因子，故有限生成且投射。任意左模同态 \(Re\rightarrow R\) 由 \(e\) 的像 \(b\in eR\) 决定，公式为 \(re\mapsto rb\)。因此若 \(Re\) 为生成元，就有某个满射 \((Re)^{(L)}\rightarrow R\)。给 \(1\) 取一个原像，这个原像仅有有限支集；结合各分量上的同态公式便得到 \(1=\sum_{i=1}^n a_i e b_i\)，所以 \(ReR=R\)。此处和式有限来自这个原像的有限支集。

反过来，若 \(1=\sum_{i=1}^n a_i e b_i\)，映射 \((Re)^n\rightarrow R\)、\((p_i)\mapsto\sum_i p_i b_i\) 的像包含一，因而满射。其截面可取 \(r\mapsto(ra_i e)_i\)，所以 \(R\) 是 \((Re)^n\) 的直和因子。\cref{b2:thm:generator-characterizations}使 \(Re\) 为生成元。

若 \(ReR=R\) 且 \(eM=0\)，对每个 \(m\in M\) 有 \(m=\sum_i a_i e b_i m=0\)，所以 \(M=0\)。若 \(ReR\ne R\)，则 \(M=R/ReR\) 为非零幺左模，而 \(eM=0\)。这就证明第三个条件与满性等价，也说明非满的取角会把某个非零模送到零。

由\cref{b2:prop:corner-endomorphism}，\(\operatorname{End}_R(Re)^{\mathrm{op}}\cong S\)。与该命题中的同态识别一样，一个 \(f:Re\rightarrow M\) 完全由 \(f(e)\) 决定，因为 \(f(re)=rf(e)\)，且 \(f(e)=ef(e)\in eM\)。因此有自然同构
\[
\operatorname{Hom}_R(Re,M)\longrightarrow eM,\qquad f\longmapsto f(e),
\]
逆为 \(m\mapsto(re\mapsto rm)\)。若 \(re=r'e\)，因 \(m=em\) 有 \(rm=r'm\)，故逆公式良定义。对 \(s\in S\)，右乘 \(s\) 是 \(Re\) 的自同态，而 Hom 中的左作用满足 \((s\cdot f)(e)=f(es)=sf(e)\)，所以此同构保留左 \(S\) 作用。这里 \(eM\) 一般不是 \(M\) 的左 \(R\) 子模，所用的标量环是 \(S\)。当 \(e\) 满时，\(Re\) 已是有限生成投射生成元，\cref{b2:thm:morita-equivalence}给出题述等价。

求值映射 \(\varepsilon_M\) 也有直接的逆公式。固定一次表达 \(1=\sum_i a_i e b_i\)，则 \(m=\sum_i a_i e b_i m\)；要使求值恢复 \(m\)，就把每一项拆成 \(a_i e\otimes e b_i m\)。由此得到加性映射
\begin{equation}\label{b2:eq:full-idempotent-evaluation-inverse}
\delta_M(m)=\sum_i a_i e\otimes e b_i m.
\end{equation}
于是 \(\varepsilon_M\delta_M(m)=\sum_i a_i e b_i m=m\)。反向复合在 \(re\otimes m\)、\(m=em\) 上得到
\[
\sum_i a_i e\otimes e b_i r e m
=\sum_i a_i e b_i r e\otimes m=re\otimes m,
\]
其中中间一步移动的平衡标量是 \(e b_i r e\in S\)。故 \(\delta_M\) 与左 \(R\) 线性映射 \(\varepsilon_M\) 互为逆，\(\delta_M\) 也左 \(R\) 线性。求值与模同态相容，因而此同构自然。另一方面，\(e(Re\otimes_SN)\) 自然识别为 \(S\otimes_SN\cong N\)：右 \(S\) 模分裂投影 \(Re\rightarrow S\)、\(p\mapsto ep\) 与其包含映射张量后仍互为投影与截面，故它的像正是 \(e(Re\otimes_SN)\)，最后使用单位张量同构。
\end{proof}

对含幺环 \(B\) 及 \(n\ge1\)，在矩阵环 \(R=M_n(B)\) 中，\(e=E_{11}\) 满，因为 \(1=\sum_iE_{i1}eE_{1i}\)。上述公式正好恢复\cref{b2:prop:matrix-morita-equivalence}。当 \(B=k\) 为域且 \(n>1\) 时，令 \(A=M_n(k)\)，则幺左 \(A\) 模 \(Ae\) 是有限生成投射生成元，却不是自由模：\(\dim_kAe=n\)，而任意非零自由左 \(A\) 模都包含一个正则模 \(A\) 的直和分量，故其 \(k\) 维数至少为 \(n^2>n\)，无论自由基有限还是无限都不可能同构于 \(Ae\)。

\ProseParagraphBreak
若 \(R=B\times C\)、\(e=(1,0)\)，且 \(C\ne0\)，则 \(ReR=B\times0\ne R\)，取角只保留第一个坐标；第二个坐标上的非零模被送到零，所以不能给出等价。

\ProseParagraphBreak
对有限维代数 \(A\)，\cref{b2:prop:basic-corner}从每种有限生成不可分解投射模的同构类型中选取一份，构造了满足 \(AeA=A\) 的幂等元，并证明 \(eAe\) 为基本代数。\cref{b2:prop:full-idempotent-morita}进一步说明 \(A\) 与这个基本代数的全部幺左模范畴等价。因此可先在基本代数上研究模，再经上述张量函子返回原代数。

\begin{proposition}[满角环的双边理想对应]\label{b2:prop:full-corner-ideal-correspondence}
设 \(R\) 为含幺环，\(e\in R\) 为满幂等元，即 \(e^2=e\) 且 \(ReR=R\)，令 \(S=eRe\)。对双边理想 \(I\subseteq R\) 与 \(K\subseteq S\)，映射
\[
I\longmapsto eIe,\qquad K\longmapsto RKR
\]
给出 \(R\) 与 \(S\) 的双边理想格之间互逆且保持包含关系的对应，其中 \(RKR\) 表示元素 \(rkt\)、\(r,t\in R\)、\(k\in K\) 的有限和组成的双边理想。对 \(R\) 的每个双边理想 \(I\)，记 \(\bar e=e+I\)，则 \(\bar e\) 是 \(R/I\) 的满幂等元，并有环同构
\[
\bar e(R/I)\bar e\cong S/eIe.
\]
因此 \(R/I\) 与 \(S/eIe\) 的全部幺左模范畴等价。
\end{proposition}
\begin{proof}
若 \(I\) 为 \(R\) 的双边理想，则 \(eIe\) 为 \(S\) 的双边理想；若 \(K\) 为 \(S\) 的双边理想，则上述有限和组成的 \(RKR\) 为 \(R\) 的双边理想。固定 \(1=\sum_{i=1}^n a_i e b_i\)。对 \(x\in I\)，两次使用这个等式得到
\[
x=\sum_{i,j}a_i e(b_i x a_j)e b_j\in R(eIe)R.
\]
反向包含 \(R(eIe)R\subseteq I\) 来自 \(I\) 的双边理想性质，故 \(R(eIe)R=I\)。

对 \(k\in K\) 有 \(k=eke\)，所以 \(K\subseteq e(RKR)e\)。反向包含可逐项检查：若 \(r,t\in R\)、\(k\in K\)，则
\[
e(rkt)e=(ere)k(ete)\in K.
\]
故 \(e(RKR)e=K\)。两种映射显然都保持包含关系，因此给出互逆的理想格对应。

商映射限制到 \(S\) 给出满环同态
\[
S\longrightarrow\bar e(R/I)\bar e,\qquad
ere\longmapsto\bar e(r+I)\bar e.
\]
其核为 \(S\cap I=eIe\)，因为 \(s\in S\cap I\) 时 \(s=ese\)。商环同构定理便给出所述角环同构。\(1=\sum_i a_i e b_i\) 在 \(R/I\) 中仍成立，故 \(\bar e\) 满；由\cref{b2:prop:full-idempotent-morita}应用于 \(R/I\) 即得模范畴等价。此论证也包含 \(I=R\) 的情形，此时两个商环均为零环。
\end{proof}
